% !TeX root = linear-codes.tex
\section{Derivation of the moment and tail estimates}
\label{sec:moment-tail-derivation}

We derive \autoref{lem:ball-fourier-estimates} from
\cite[Corollary 1.4 and Theorem 1.5]{KirshnerS2021}.
Put $\rho=w/n$.  The third-moment specialization of their moment bound
says that every polynomial $f$ of degree at most $w$ satisfies
$$
\E|f|^3
\le
(\E f^2)^{3/2}
2^{n(\entb{\rho}/2-\gamma(\rho))},
$$
where
$$
\begin{aligned}
\gamma(\rho)\defeq{}&
-\frac{\rho}{2}\lg \rho-2(1-\rho)\lg(1-\rho)-1\\
&+\left(1-\frac{3\rho}{2}\right)\lg(2-3\rho).
\end{aligned}
$$
Indeed, the auxiliary parameter in their moment exponent satisfies
$$
1-2u=\sqrt{\frac{\rho}{2-3\rho}},
$$
and substitution gives the displayed expression.  Elementary
calculus gives
$$
\gamma'(\rho)
=\frac12\lg\frac{(1-\rho)^4}{\rho(2-3\rho)^3}.
$$
It follows that $\gamma$ is strictly positive on $(0,1/2)$ and
$$
\gamma(\rho)=\frac{\rho}{2}\lg\frac1\rho+O(\rho)
\qquad(\rho\longrightarrow0).
$$
Consequently, for every fixed $\eta>0$, there is a constant
$c_\eta>0$ such that
$$
\gamma(\rho)\ge c_\eta \rho\lg\frac1\rho
\qquad
\left(0<\rho\le\frac12-\eta\right).
$$
As a function of the signs
$((-1)^{x_1},\ldots,(-1)^{x_n})$, the function
$\widehat\mu(x)$ is a multilinear polynomial of degree at most
$w$.  Moreover, Parseval's identity and Stirling's formula give
$$
\E_x\widehat\mu(x)^2=\frac1V,
\qquad
\lg V=n\entb{\rho}+O(\lg n).
$$
Applying the preceding moment estimate to $\widehat\mu$ gives
$$
\begin{aligned}
Vn^2\E_x|\widehat\mu(x)|^3
&\le
n^2V^{-1/2}
2^{n\entb{\rho}/2-c_\eta w\lg(n/w)}\\
&=
2^{-c_\eta w\lg(n/w)+O(\lg n)}
=O(1/n),
\end{aligned}
$$
where the last equality follows from $w\gg\ln n$.
This proves the moment estimate in
\autoref{lem:ball-fourier-estimates}.

We next derive the tail estimate.  For
$$
0\le y\le\frac12-\sqrt{\rho(1-\rho)}
$$
define
$$
T(\rho,y)
\defeq
\entb{\rho}+
\int_0^y
\lg\left(
\frac{1-2\rho+\sqrt{(1-2\rho)^2-4z(1-z)}}{2(1-z)}
\right)\,dz.
$$
Kirshner and Samorodnitsky's tail bound says that, when $ny$ is an
integer,
$$
\Pr\left[
|f|\ge
\sqrt{\E f^2}\,
2^{n(T(\rho,y)-\entb{\rho}/2)}
\right]
\le
2^{n(\entb{y}-1)}.
$$
The expression inside the logarithm defining $T$ is at most
$1-2\rho$.  Therefore
$$
\entb{\rho}-T(\rho,y)
\ge y\lg\frac1{1-2\rho}
\ge\frac{2\rho y}{\ln2}.
$$

Take
$$
y=\frac1n\left\lceil\frac{n\lg n}{w}\right\rceil.
$$
Our assumption $w\gg\ln n$ gives $y\longrightarrow0$, so the
permitted range for $y$ holds for all sufficiently large $n$.
The standard estimate
$$
V\ge\binom nw\ge\frac{2^{n\entb{\rho}}}{n+1}
$$
and the preceding bound on $T$ give
$$
\begin{aligned}
V^{-1/2}2^{n(T(\rho,y)-\entb{\rho}/2)}
&\le
\sqrt{n+1}\,2^{n(T(\rho,y)-\entb{\rho})}\\
&\le
\sqrt{n+1}\,2^{-2wy/\ln2}\\
&\le
\sqrt{n+1}\,n^{-2/\ln2}\\
&\le\frac1n
\end{aligned}
$$
for all sufficiently large $n$.  Hence the threshold in the
tail bound is at most $1/n$.  Moreover,
$$
y\le\frac{\lg n}{w}+\frac1n
\le\frac{2\lg n}{w}.
$$
Using $\entb{y}\le y\lg(e/y)$, we conclude that
$$
\begin{aligned}
\Pr_x\sparen{|\widehat\mu(x)|>1/n}
&\le
2^{-n+n\entb{y}}\\
&\le
2^{-n+\frac{2n\lg n}{w}
\lg\left(\frac{ew}{\lg n}\right)}.
\end{aligned}
$$
This proves the tail estimate and completes the derivation of
\autoref{lem:ball-fourier-estimates}.
